(1) 由于内部工质经过一循环不积累熵，吸收热量与放出热量的熵量相同，即 \(K(T_h-X)/X=K(Y-T_c)/Y\)。整理为 \(T_h/X+T_c/Y=2\)。输出功率为 \(\dot W=\dot Q_h-\dot Q_c=K(T_h+T_c-X-Y)\)，必须同时满足 \(T_h>X>Y>T_c\) 才是热机。

(2) 令 \(r=Y/X\)，则由约束 \(X=(T_h+T_c/r)/2\)、\(Y=(rT_h+T_c)/2\)。输出功率为 \(\dot W(r)=K[T_h(1-r)+T_c(1-1/r)]/2\)，效率为 \(\eta=1-Y/X=1-r\)。可行的热机工作区间为 \(T_c/T_h<r<1\)，其中两端功率为零。

(3) 求导 \(d\dot W/dr=K[-T_h+T_c/r^2]/2\)，二阶导数为 \(-KT_c/r^3<0\)，唯一极大点 \(r_*=\sqrt{T_c/T_h}\)。故 \(\eta_*=1-\sqrt{T_c/T_h}\)，\(\dot W_{\max}=K(\sqrt{T_h}-\sqrt{T_c})^2/2\)。因为平方根大于原始 \(T_c/T_h\)，最大功率效率低于可逆两热库卡诺效率。

(4) 设 \(T_h=T+\Delta/2\)、\(T_c=T-\Delta/2\)，\(|\Delta|\ll T\) 时 \(\sqrt{T_h}-\sqrt{T_c}\simeq\Delta/(2\sqrt T)\)，所以 \(\dot W_{\max}\simeq K\Delta^2/(8T)\)。形式上 \(K\to\infty\) 使最大功率发散，只是无限换热能力、忽略有限面积和材料限制的理想化表现；真实内外热阻及动力学耗散使功率有上限。
